\section{Sistemas de datos: calidad, licencias y la triple restricción de la «contaminación de datos»}

\subsection{La mezcla de datos y los datos sintéticos se vuelven capacidades habituales}

La práctica de ingeniería de datos mencionada en la conversación es muy pragmática: primero se hacen experimentos de mezcla con muestras pequeñas y después se ajusta la proporción óptima según la evaluación objetivo. Este proceso refleja el principio de «datos orientados a la tarea» y no el de «entre más datos, mejor». Al mismo tiempo, el OCR y las muestras sintéticas de alta calidad permiten que los tokens disponibles sigan creciendo.

\begin{importantbox}{La forma correcta de usar datos sintéticos}
El valor real de los datos sintéticos depende de que exista un «filtrado verificable»:
\begin{itemize}[leftmargin=1.5em]
    \item Problemas matemáticos: verificación programática;
    \item Problemas de código: pruebas unitarias y verificación por ejecución;
    \item Preguntas y respuestas complejas: consistencia cruzada entre varios modelos más revisión humana por muestreo.
\end{itemize}
Esto significa que la «generación» es apenas el inicio; lo determinante está en el filtrado, el etiquetado y el control de versiones posteriores.
\end{importantbox}

\subsection{Límites legales y soberanía de datos}

La entrevista aborda un riesgo legal real: si el corpus de entrenamiento cuenta con autorización, si el rastreo histórico cumple la normativa y cómo se rastrea la responsabilidad de derechos de autor. Este problema no se resolverá por sí solo con el avance tecnológico. Cuanto más se avance, más se convertirá la soberanía de datos en una ventaja competitiva duradera para las organizaciones: quien pueda obtener legalmente datos privados de alta calidad tendrá más oportunidades de desarrollar capacidades que no sean homogéneas frente a las demás.

\begin{warningbox}{«Si se puede obtener, se puede entrenar con eso» es una idea de alto riesgo}
En un entorno regulatorio de alta presión, la cadena de comprobación del origen de los datos se parecerá cada vez más a una auditoría financiera. Un corpus sin origen rastreable puede acelerar el entrenamiento a corto plazo, pero a largo plazo puede provocar indemnizaciones cuantiosas, el retiro de modelos y daños a la reputación de la marca, riesgos que superan por mucho el beneficio de un solo entrenamiento.
\end{warningbox}

\subsection{La contaminación de datos de los LLM y la carga sobre los mantenedores de código abierto}

A medida que aumenta la proporción de contenido generado por IA en la red, el corpus de entrenamiento comienza a presentar el riesgo circular de que «un modelo aprenda de otro modelo». La observación de Sebastian sobre la comunidad de mantenimiento de código abierto es muy concreta: una gran cantidad de pull requests de baja calidad asistidos por IA está consumiendo la energía de los mantenedores. Lo determinante aquí no es la «generación por IA» en sí, sino si existe una capa de verificación humana.

\begin{knowledgebox}{Un criterio sencillo para juzgar qué son «datos utilizables»}
Si un fragmento de contenido no puede ejecutarse, contrastarse, reproducirse ni verificarse de forma cruzada, difícilmente podrá convertirse en una muestra de entrenamiento de alto valor. El corpus de código es relativamente fácil de procesar porque puede ejecutarse; los textos de explicación conceptual y de comentario son más difíciles de evaluar en cuanto a calidad y requieren canales de evaluación adicionales.
\end{knowledgebox}

\subsection{Resumen de la sección}

Para 2026, la ingeniería de datos ha entrado en la etapa de la «gestión de la calidad de datos»: la estrategia de mezcla, el cumplimiento legal y el control de la contaminación tienen la misma importancia; ninguno puede faltar.

